\section*{Bab \ref{ch:g11:cell-cycle-mitosis} --- Daur Sel dan Mitosis}

\begin{solution}{exo:g11:cell-cycle-mitosis:1}
$\mathrm{G_1}$: pertumbuhan; S: \omterm{def:g11:dna-replication:replication}{replikasi DNA}; $\mathrm{G_2}$:
persiapan pembelahan; M: \omterm{def:g11:cell-cycle-mitosis:cycle}{mitosis} (pembelahan \omterm{def:g10:cells-common-unit:organelle}{inti selnya}) dan
sitokinesis (pembelahan \omterm{def:g10:cells-common-unit:cell}{sitoplasmanya}).
\end{solution}

\begin{solution}{exo:g11:cell-cycle-mitosis:2}
\omterm{def:g11:cell-cycle-mitosis:chromosome}{Kromatid} adalah salah satu dari dua salinan serupa sebuah \omterm{def:g11:cell-cycle-mitosis:chromosome}{kromosom} yang
telah direplikasi; \omterm{def:g11:cell-cycle-mitosis:chromosome}{sentromer} adalah daerah tempat keduanya tertahan
bersama. Sebuah \omterm{def:g11:cell-cycle-mitosis:chromosome}{kromosom} punya dua \omterm{def:g11:cell-cycle-mitosis:chromosome}{kromatid} sejak akhir S sampai anafase.
\end{solution}

\begin{solution}{exo:g11:cell-cycle-mitosis:3}
Profase: \omterm{def:g11:cell-cycle-mitosis:chromosome}{kromosomnya} memadat, selubung intinya terurai. Metafase:
\omterm{def:g11:cell-cycle-mitosis:chromosome}{kromosomnya} berbaris pada khatulistiwanya. Anafase: \omterm{def:g11:cell-cycle-mitosis:chromosome}{kromatidnya} memisah
dan bergerak ke kutubnya. Telofase: selubung intinya terbentuk lagi,
lalu sitokinesis.
\end{solution}

\begin{solution}{exo:g11:cell-cycle-mitosis:4}
Dalam fase S, setengah jalan replikasinya.
\end{solution}

\begin{solution}{exo:g11:cell-cycle-mitosis:5}
46 \omterm{def:g11:cell-cycle-mitosis:chromosome}{kromosom}, 92 \omterm{def:g11:cell-cycle-mitosis:chromosome}{kromatid}.
\end{solution}

\begin{solution}{exo:g11:cell-cycle-mitosis:6}
$\mathrm{G_1}$ \qty{11}{h}, S \qty{8}{h}, $\mathrm{G_2}$ \qty{4}{h},
M \qty{1}{h}. \omterm{def:g10:cells-common-unit:cell}{Sel} dengan $2q$ berada dalam $\mathrm{G_2}$ atau dalam
\omterm{def:g11:cell-cycle-mitosis:cycle}{mitosis} sebelum akhir anafasenya.
\end{solution}

\begin{solution}{exo:g11:cell-cycle-mitosis:7}
Metafase: 40 \omterm{def:g11:cell-cycle-mitosis:chromosome}{kromosom}, 80 \omterm{def:g11:cell-cycle-mitosis:chromosome}{kromatid}. Anafase: 80 \omterm{def:g11:cell-cycle-mitosis:chromosome}{kromosom} (berkromatid
tunggal), 40 ke tiap kutubnya. \omterm{def:g10:cells-common-unit:cell}{Sel} anaknya: 40 \omterm{def:g11:cell-cycle-mitosis:chromosome}{kromosom}, 40 \omterm{def:g11:cell-cycle-mitosis:chromosome}{kromatid}.
\end{solution}

\begin{solution}{exo:g11:cell-cycle-mitosis:8}
Indeks \omterm{def:g11:cell-cycle-mitosis:cycle}{mitosisnya} $80/800 = 10\%$; \omterm{def:g11:cell-cycle-mitosis:cycle}{mitosisnya} \qty{2.4}{h}. Profase
$48/800 \times 24 = \qty{1.44}{h}$; metafase \qty{0.48}{h} (29 menit);
anafase \qty{0.18}{h} (11 menit); telofase \qty{0.3}{h} (18 menit).
\end{solution}

\begin{solution}{exo:g11:cell-cycle-mitosis:9}
Tanpa gelendong, \omterm{def:g11:cell-cycle-mitosis:chromosome}{kromatidnya} tidak dapat ditarik terpisah: \omterm{def:g10:cells-common-unit:cell}{selnya}
tersangkut dalam metafase dengan \omterm{def:g11:cell-cycle-mitosis:chromosome}{kromosomnya} yang memadat dan
sepenuhnya terlihat, masing-masing dengan dua \omterm{def:g11:cell-cycle-mitosis:chromosome}{kromatid}. Itulah tepat
keadaan ketika \omterm{def:g11:cell-cycle-mitosis:chromosome}{kromosom} dipotret untuk sebuah \omterm{def:g11:cell-cycle-mitosis:chromosome}{kariotipe}.
\end{solution}

\begin{solution}{exo:g11:cell-cycle-mitosis:10}
Dua meter benang yang kusut tidak dapat dipilah; setelah memadat
menjadi batang sepanjang beberapa mikrometer, ke-46 \omterm{def:g11:cell-cycle-mitosis:chromosome}{kromosomnya} dapat
dibariskan dan ditarik terpisah tanpa putus atau tersimpul. \omterm{def:g11:cell-cycle-mitosis:chromosome}{Kromosomnya}
terurai sesudahnya karena \omterm{def:g10:universal-dna:gene}{gen} hanya dapat dibaca dari bentuk yang longgar.
\end{solution}

\begin{solution}{exo:g11:cell-cycle-mitosis:11}
Sepuluh pembelahan: $2^{10} = 1024$ \omterm{def:g10:cells-common-unit:cell}{sel}. Tidak: embrio awalnya tidak
tumbuh di antara pembelahannya, dan daurnya hampir seluruhnya berupa S
dan M, dengan $\mathrm{G_1}$ yang nyaris tidak ada.
\end{solution}

\begin{solution}{exo:g11:cell-cycle-mitosis:12}
Satu anaknya punya 47 \omterm{def:g11:cell-cycle-mitosis:chromosome}{kromosom} (satu berlebih), yang lain 45 (satu
hilang); \omterm{def:g10:universal-dna:information}{DNA-nya} $q$ ditambah atau dikurangi jatah satu \omterm{def:g11:cell-cycle-mitosis:chromosome}{kromosom}. Dalam
\omterm{def:g10:cells-common-unit:cell}{sel} tubuh, kesalahan itu mengenai satu \omterm{def:g10:cells-common-unit:cell}{sel} di antara triliunan, yang
lazimnya mati atau disingkirkan; tetapi jika \omterm{def:g10:cells-common-unit:cell}{selnya} bertahan dan
membelah, dan \omterm{def:g11:cell-cycle-mitosis:chromosome}{kromosom} yang hilang atau bertambah itu mengacaukan
pengaturan pertumbuhannya, sebuah tumor dapat mulai.
\end{solution}

\begin{solution}{exo:g11:cell-cycle-mitosis:13}
Sekitar $\num{2e11}$ pembelahan tiap hari, yaitu $\num{2e11}$ \omterm{def:g10:cells-common-unit:cell}{sel} yang
berdaur pada suatu saat (masing-masing memberi satu \omterm{def:g10:cells-common-unit:cell}{sel} baru tiap
hari). Menghalangi \omterm{def:g11:cell-cycle-mitosis:cycle}{mitosis} menghentikan pasokannya; \omterm{def:g10:cells-common-unit:cell}{sel} merah hidup 120
hari tetapi sebagian \omterm{def:g10:cells-common-unit:cell}{sel} putih hanya beberapa hari, sehingga dalam
sepekan cacah \omterm{def:g10:cells-common-unit:cell}{sel} putihnya runtuh dan infeksi menyusul --- itulah efek
samping yang khas obat antikanker.
\end{solution}

\begin{solution}{exo:g11:cell-cycle-mitosis:14}
Dalam interfase \omterm{def:g11:cell-cycle-mitosis:chromosome}{kromosomnya} terurai dan tak terlihat; dalam anafase
\omterm{def:g11:cell-cycle-mitosis:chromosome}{kromosomnya} sedang bergerak dan bercampur. Dalam metafase \omterm{def:g11:cell-cycle-mitosis:chromosome}{kromosomnya}
sepenuhnya memadat, terpisah, masih terbuat dari dua \omterm{def:g11:cell-cycle-mitosis:chromosome}{kromatid} dan
terletak dalam satu bidang: masing-masing dapat dikenali dari ukuran
dan bentuknya.
\end{solution}

\begin{solution}{exo:g11:cell-cycle-mitosis:15}
\omterm{def:g10:cells-common-unit:cell}{Sel} hati yang beristirahat dalam $\mathrm{G_1}$ masuk lagi ke daurnya,
melewati S, $\mathrm{G_2}$ dan \omterm{def:g11:cell-cycle-mitosis:cycle}{mitosis} berulang kali sampai massanya
pulih, lalu kembali beristirahat. Pada hari ke-0 hampir seluruh \omterm{def:g10:cells-common-unit:cell}{selnya}
memuat $q$; pada hari ke-3 banyak yang memuat antara $q$ dan $2q$
(dalam S) atau $2q$ (dalam $\mathrm{G_2}$, \omterm{def:g11:cell-cycle-mitosis:cycle}{mitosis}).
\end{solution}

\begin{solution}{pb:g11:cell-cycle-mitosis:1}
\textbf{1.} $100/1000 = 10\%$.

\textbf{2.} $0.10 \times 20 = \qty{2}{h}$.

\textbf{3.} Profase $60/1000 \times 20 = \qty{1.2}{h}$ (72 menit);
metafase $0.36$ jam (22 menit); anafase $0.16$ jam (sekitar 10 menit);
telofase $0.28$ jam (17 menit).

\textbf{4.} Anafase. \omterm{def:g11:cell-cycle-mitosis:chromosome}{Kromatidnya} sudah memadat dan melekat; gerakannya
berupa tarikan sejauh beberapa mikrometer oleh serat yang memendek
dalam hitungan menit.

\textbf{5.} Tahap yang langka terlihat pada sedikit \omterm{def:g10:cells-common-unit:cell}{sel}, dan hitungan
yang kecil memberi ralat yang besar; lagi pula hanya zona
pertumbuhannya yang berdaur --- di tempat lain indeksnya akan nol dan
mengencerkan angkanya.

\textbf{6.} 16 \omterm{def:g11:cell-cycle-mitosis:chromosome}{kromosom}, 32 \omterm{def:g11:cell-cycle-mitosis:chromosome}{kromatid}.

\textbf{7.} 32 \omterm{def:g11:cell-cycle-mitosis:chromosome}{kromosom} yang bergerak, 16 ke tiap kutubnya.

\textbf{8.} Profase $2q$; anafase $2q$ dalam seluruh \omterm{def:g10:cells-common-unit:cell}{selnya}; telofase
$q$ tiap \omterm{def:g10:cells-common-unit:organelle}{inti selnya}.

\textbf{9.} S adalah $7/19$ interfasenya: sekitar $900 \times 7/19 \approx
330$ \omterm{def:g10:cells-common-unit:cell}{sel}.

\textbf{10.} Anafase: \omterm{def:g11:cell-cycle-mitosis:chromosome}{sentromernya} baru saja membelah dan ke-32
\omterm{def:g11:cell-cycle-mitosis:chromosome}{kromatid} dari 16 \omterm{def:g11:cell-cycle-mitosis:chromosome}{kromosomnya} kini menjadi 32 \omterm{def:g11:cell-cycle-mitosis:chromosome}{kromosom} yang sedang menuju
kutubnya.

\textbf{11.} $\num{20000}/20 = 1000$ \omterm{def:g10:cells-common-unit:cell}{sel} baru tiap jam.

\textbf{12.} 1000 \omterm{def:g10:cells-common-unit:cell}{sel} tiap jam pada seluruh zonanya; dalam satu deret,
zonanya barangkali selebar seratus \omterm{def:g10:cells-common-unit:cell}{sel} pada tiap arah, sehingga kira-kira
$\num{24000}/(100 \times 100) \approx 2$--3 \omterm{def:g10:cells-common-unit:cell}{sel} tiap deret tiap hari,
masing-masing \qty{100}{\micro m}: beberapa persepuluh milimeter.
Angka \qty{1}{mm} yang teramati berarti zonanya lebih sempit daripada
yang dianggap atau \omterm{def:g10:cells-common-unit:cell}{selnya} memanjang lebih jauh; orde besarnya benar.

\textbf{13.} Satu titik awal menyalin $2 \times 50 \times 7 \times 3600 =
\num{2.5e6}$ pasangan \omterm{def:g10:universal-dna:nucleotide}{basa}; $\num{1.6e10}/\num{2.5e6} \approx 6400$
titik awal sekurang-kurangnya.

\textbf{14.} Tiap tahapnya berlangsung lebih lama, sehingga lebih
sedikit \omterm{def:g10:cells-common-unit:cell}{sel} yang merampungkan daurnya tiap hari; jika semua fasenya
melambat sama rata, bagiannya, dan dengan demikian indeks \omterm{def:g11:cell-cycle-mitosis:cycle}{mitosisnya},
tidak berubah.

\textbf{15.} \omterm{def:g10:cells-common-unit:cell}{Sel} yang sudah melewati S akan merampungkan
$\mathrm{G_2}$ dan \omterm{def:g11:cell-cycle-mitosis:cycle}{mitosisnya} pada jam-jam pertamanya; setelah sehari
tidak ada \omterm{def:g10:cells-common-unit:cell}{sel} yang dapat mencapai \omterm{def:g11:cell-cycle-mitosis:cycle}{mitosis}, sehingga profase sampai
telofase telah lenyap dan hanya \omterm{def:g10:cells-common-unit:cell}{sel} interfase yang tersisa, tersangkut
dalam $\mathrm{G_1}$ atau pada awal S.

\textbf{16.} $4/1000 \times 20 = \qty{0.08}{h}$, sekitar 5 menit,
berhadapan dengan 10. Tidak mengejutkan: dengan 4 atau 8 \omterm{def:g10:cells-common-unit:cell}{sel},
hitungannya dikuasai kebetulan.

\textbf{17.} $12/2000 \times 20 = \qty{0.12}{h}$, sekitar 7 menit.

\textbf{18.} Manusia: $0.04 \times 24 \approx \qty{1}{h}$; bawang
putihnya \qty{2}{h}. Ordenya serupa, \omterm{def:g10:cells-common-unit:cell}{sel} tumbuhannya agak lebih lambat.

\textbf{19.} \omterm{def:g10:cells-common-unit:cell}{Sel} tumor berdaur terus-menerus sedangkan sebagian besar
\omterm{def:g10:cells-common-unit:cell}{sel} normal beristirahat dalam $\mathrm{G_1}$, sehingga bagian yang lebih
besar berada dalam \omterm{def:g11:cell-cycle-mitosis:cycle}{mitosis} pada suatu saat. Karena itu obat penghalang
gelendong mengenai \omterm{def:g10:cells-common-unit:cell}{sel} tumor lebih sering daripada \omterm{def:g10:cells-common-unit:cell}{sel} normal ---
tetapi juga mengenai jaringan normal yang cepat membelah.

\textbf{20.} Daurnya 20 jam: $\mathrm{G_1}$ 8, S 7, $\mathrm{G_2}$ 4,
\omterm{def:g11:cell-cycle-mitosis:cycle}{mitosisnya} sekitar 2 (profase 72 menit, metafase 22, anafase 7--10,
telofase 17). Anafase, yang terpendek, adalah saat \omterm{def:g11:cell-cycle-mitosis:chromosome}{sentromernya}
membelah dan tiap kutubnya menerima satu \omterm{def:g11:cell-cycle-mitosis:chromosome}{kromatid} dari masing-masing
keenam belas \omterm{def:g11:cell-cycle-mitosis:chromosome}{kromosomnya}.
\end{solution}
